% Generated draft: requires mathematical and source review.
\documentclass[11pt]{article}
\usepackage[margin=25mm]{geometry}
\usepackage{fontspec}
\setmainfont{Latin Modern Roman}
\usepackage{amsmath,amssymb,mathtools,unicode-math}
\setmathfont{Latin Modern Math}
\usepackage{enumitem,needspace,xurl,hyperref}
\setlength{\parindent}{0pt}
\setlength{\parskip}{5pt}
\setlength{\emergencystretch}{3em}
\newcommand{\N}{\mathbb{N}}
\newcommand{\Z}{\mathbb{Z}}
\newcommand{\Q}{\mathbb{Q}}
\newcommand{\R}{\mathbb{R}}
\newcommand{\C}{\mathbb{C}}
\newcommand{\F}{\mathbb{F}}
\newcommand{\A}{\mathbb{A}}
\newcommand{\PP}{\mathbb{P}}
\newcommand{\E}{\mathbb{E}}
\newcommand{\eps}{\varepsilon}
\newcommand{\dd}{\,\mathrm{d}}
\begin{document}
\section*{UnsolvedMath ID: 4800015}
\textbf{Source problem code:} catalogue entry
\subsection{Definitions and mathematical statement}
The source gives integer polynomials $p_1,\ldots,p_\ell\in\Z[t]$ and asks for characteristic factors $\mathcal Z_{d,T_i}$ for averages referenced only as ``E:Multies.'' Under the standard commuting-transformation interpretation, these averages are
\[
 A_N(f_1,\ldots,f_\ell)=\frac1N\sum_{n=1}^{N}
       \prod_{i=1}^{\ell}f_i\circ T_i^{p_i(n)}
\]
on a probability space $(X,\mathcal B,\mu)$ with commuting invertible measure-preserving transformations $T_1,\ldots,T_\ell$ and bounded measurable $f_i$. This formula and ambient hypotheses must be verified against the missing reference; they are not present explicitly in the excerpt.

A usual meaning of pairwise polynomial independence in this setting is independence modulo constants: for $i\ne j$ and $(a,b)\in\Q^2\setminus\{(0,0)\}$, the polynomial $ap_i+bp_j$ is nonconstant. Ordinary linear independence differs if constants are not normalized. The source does not give the convention or a condition $p_i(0)=0$.

For a transformation $T$, $\mathcal Z_{d,T}$ denotes the Host--Kra factor of order $d$, with the customary indexing characterized by
\[
 \E(f\mid\mathcal Z_{d,T})=0\quad\Longleftrightarrow\quad
 \|f\|_{d+1,T}=0
\]
for the Host--Kra seminorm. A tuple of these factors is characteristic in $L^2$ if
\[
 \left\|A_N(f_1,\ldots,f_\ell)-
 A_N\bigl(\E(f_1\mid\mathcal Z_{d,T_1}),\ldots,
          \E(f_\ell\mid\mathcal Z_{d,T_\ell})\bigr)\right\|_{L^2(\mu)}\to0.
\]
Equivalently, a sufficient vanishing formulation requires $A_N\to0$ in $L^2$ whenever one of these conditional expectations is zero and the other functions remain arbitrary bounded functions.

\textbf{Problem (Nikos Frantzikinakis).} For every pairwise independent integer polynomial family in the intended setting, show that some finite common order $d\in\N$ makes $\mathcal Z_{d,T_1},\ldots,\mathcal Z_{d,T_\ell}$ characteristic. The source does not specify whether $d$ depends only on the polynomial family or also on the system; no stronger uniform dependence is silently imposed. Nonconstant-polynomial and $\ell$ conventions also require context, especially when $\ell=1$ makes pairwise independence vacuous.

The summary reports a solution by F.--Kuca (\emph{Inventiones}, 2025), also claiming an associated joint-ergodicity conclusion without its hypotheses. The background separately says the author's November 2025 progress page did not record a resolution of this numbered item. These evidence statements are retained without asserting that either authenticates the exact missing-average formulation.
\subsection{Short English statement}
Find a finite common Host--Kra order giving characteristic factors for the multiple averages associated with pairwise independent integer polynomials and commuting transformations. The standard average is displayed conditionally because its source reference is missing. The registry reports a 2025 solution, while its background retains an unresolved progress-page note.
\subsection{Sources}
\begin{itemize}
\item Local UnsolvedMath registry, permanent ID 4800015, source code catalogue entry, ``Multiple ergodic averages --- Problem 15''; proposer: Nikos Frantzikinakis, proposed year 2016. Supplied status solved, classification SOLVED-IN-LITERATURE, updated 2026-08-06.
\item Frantzikinakis, \emph{Some open problems on multiple ergodic averages} (source-list date 2016), Problem 15, \url{https://arxiv.org/abs/1103.3808}. Reported source-TeX extraction and access: 2026-07-29. The record says a November 2025 author progress page did not record resolution of this numbered problem, without supplying that page's URL. Rights: NEEDS\_REVIEW; public access verified but redistribution permission not established.
\item The supplied research summary cites ``F.--Kuca (Inventiones 2025)'' without a title, theorem number, or URL, and reports characteristic-factor and joint-ergodicity conclusions. The reference E:Multies is unresolved in the extracted statement.
\end{itemize}
\end{document}
